\section{Concrete Security Summary}
\label{sec:concrete-security}

\Cref{sec:aead-security,sec:root-tag-hash-security} separate the security
analysis into AEAD bounds and root tag hash and committing bounds. This section
instantiates those bounds at the
implemented \TW{} parameters and records the single concrete work cap used
throughout the paper.

\subsection{Concrete TW128 Bounds}
\label{sec:concrete}

The implemented parameters of \Cref{sec:params} fix
\[
  k = c = \nu = \tauroot = \tauleaf = 256.
\]
Substituting into
\Cref{cor:ae,thm:mu-ind,thm:int-rup,thm:coll,thm:pre,cor:second-preimage-rs04,cor:cmtds4,thm:epre}
gives the following closed-form bounds.

\paragraph{All-in-one AE.}
\[
  \Advae_\TW(\mathcal{A})
  \;\leq\;
  \frac{\Ntot(\Ntot + 1)}{2^{257}}
  + \frac{\Nprim(\Nprim + 1)}{2^{257}}
  + \frac{\Nprim + \Qv + \Nleaf}{2^{256}}.
\]

\paragraph{Multi-user indistinguishability.}
\[
  \Advmuind_\TW(\mathcal{A})
  \;\leq\;
  \frac{\Ntot(\Ntot + 1)}{2^{257}}
  + \frac{\Nprim(\Nprim + 1)}{2^{257}}
  + \frac{D \Nprim + \Qv + \Nleaf}{2^{256}}
  + \frac{D(D - 1)}{2^{257}}.
\]

\paragraph{\textsf{INT-RUP} integrity.}
\[
  \AdvintRUP_\TW(\mathcal{A})
  \;\leq\;
  \frac{\Ntot(\Ntot + 1)}{2^{257}}
  + \frac{\Nprim + \Qv + \Nleaf}{2^{256}}.
\]

\paragraph{$s$-way collision resistance of $\HTW$.}
For every $s \geq 2$,
\[
  \AdvColl_\TW(\mathcal{A})
  \;\leq\;
  \frac{\Ntot(\Ntot + 1)}{2^{257}}
  + \frac{\binom{\Nprim+\Nleaf}{2}}{2^{256}}
  + \frac{\binom{\Nprim + s}{s}}{2^{256 (s - 1)}}.
\]
For $s=2$, this is the ordinary two-input collision bound:
\[
\begin{aligned}
  \AdvColl_\TW(\mathcal{A})
\leq\;&
  \frac{\Ntot(\Ntot + 1)}{2^{257}}
  + \frac{(\Nprim+\Nleaf)(\Nprim+\Nleaf-1)}{2^{257}} \\
& + \frac{(\Nprim + 1)(\Nprim + 2)}{2^{257}} .
\end{aligned}
\]

\paragraph{Second-preimage resistance of $\HTW$.}
By \Cref{cor:second-preimage-rs04}, the standard and everywhere
\cite{RS04} second-preimage advantages of $\HTW$ are bounded by the same
$s=2$ collision expression above.

\paragraph{Preimage resistance of $\HTW$.}
For fixed byte lengths $a,\mu$, let
$m_{a,\mu}=k+\nu+8a+8\mu$. By \Cref{thm:pre},
\[
\begin{aligned}
  \AdvPreSlice{a}{\mu}_\TW(\mathcal{A})
\leq\;&
  \frac{\Ntot(\Ntot + 1)}{2^{257}}
  + \frac{(\Nprim+\Nleaf)(\Nprim+\Nleaf-1)}{2^{257}} \\
& + \frac{\Nprim + 1}{2^{256}}
  + 2^{256 - m_{a,\mu}} .
\end{aligned}
\]
Under the implemented parameters $k = \nu = 256$, the slice term is
$2^{-256-8a-8\mu}$.

\paragraph{$s$-way CMTDs-4 committing security.}
For every $s \geq 2$, the all-bodies-equal special case drops the
known-key leaf tag collision term:
\[
  \Advcmtds{4}_\TW(\mathcal{A})
  \;\leq\;
  \frac{\Ntot(\Ntot + 1)}{2^{257}}
  + \frac{\binom{\Nprim + s}{s}}{2^{256 (s - 1)}}.
\]

\begin{remark}[Two-Opening CMTDs-4 Specialization]
\label{rem:cmtd4-s2}
Setting $s = 2$ recovers the ordinary two-opening
CMTDs-4 statement of \cite{BH22}:
\[
  \Advcmtds{4}_\TW(\mathcal{A})
  \;\leq\;
  \frac{\Ntot(\Ntot + 1)}{2^{257}}
  + \frac{\binom{\Nprim + 2}{2}}{2^{256}}
  \;=\;
  \frac{\Ntot(\Ntot + 1) + (\Nprim + 1)(\Nprim + 2)}{2^{257}}.
\]
\end{remark}

\paragraph{Everywhere preimage resistance of $\HTW$.}
\[
  \AdvePre_\TW(\mathcal{A})
  \;\leq\;
  \frac{\Ntot(\Ntot + 1)}{2^{257}}
  + \frac{\Nprim + 1}{2^{256}}.
\]

\subsection{Concrete Security Claim}
\label{sec:concrete-claim}

\begin{corollary}[Concrete \TW{} Security]
\label{cor:concrete-security}
Under the implemented parameters $k = c = \nu = \tauroot = \tauleaf = 256$, every
adversary with work cap $\Ntot \leq 2^{63}$---and, in the multi-user case,
user cap $D \leq 2^{63}$---has advantage at most $2^{-128}$ in each of
\TW{}'s notions: multi-user indistinguishability (\Cref{thm:mu-ind}),
all-in-one AE (\Cref{cor:ae}), \textsf{INT-RUP} integrity
(\Cref{thm:int-rup}), the hash security of
$\HTW$---$s$-way collision resistance for every $s \geq 2$ in the sense
of \cite{BH22}, with the \cite{RS04} notions of standard and everywhere
second-preimage resistance inherited from collision resistance on every
fixed input slice, standard preimage resistance on every fixed input
slice, and everywhere preimage resistance
(\Cref{thm:coll,cor:second-preimage-rs04,thm:pre,thm:epre})---and the \cite{BH22}
CMTDs-4 committing bound it yields (\Cref{cor:cmtds4}).
\end{corollary}

\begin{proof}
The resource consequences of \Cref{sec:resources} give
$\Nprim \leq \Ntot$, $\Qv \leq \Ntot$, $\Nleaf \leq \Ntot$, and
$\Nprim+\Nleaf \leq \Ntot$. The user cap $D$ remains independent, since
created-but-idle users need not contribute to $N$.

At $\Ntot \leq 2^{63}$, each flat-sponge term---at $\Ntot$ or at
$\Nprim$---and each birthday-scale root or leaf term at $s=2$ is at most
$2^{-131} + 2^{-192}$, while each
linear root-preimage, key-guessing, verification, or leaf tag term with no
factor $D$ is at most a $2^{-193}$-scale term.
The fixed-slice Pre degradation is at most $2^{-256}$, since
$m_{a,\mu} \geq k+\nu = 512$.
Thus the largest non-multi-user bound is the $s=2$ collision bound, which is below
$4 \cdot 2^{-131} = 2^{-129}$; larger $s$ only decreases the
root-collision term. The all-in-one AE bound carries its two flat-sponge
terms as its only birthday-scale terms and lies below $2^{-129}$, and the
\textsf{INT-RUP} bound carries $\Lift_c(\Ntot)$ alone. In the multi-user
case, using $D \leq 2^{63}$,
\[
  \Advmuind_\TW(\mathcal{A})
  < \bigl(2^{-130} + 2^{-193}\bigr) + 2^{-130} + 2^{-192} + 2^{-131}
  < 2^{-128}:
\]
the two flat-sponge terms sum to at most
$2 \cdot (2^{-131} + 2^{-194}) = 2^{-130} + 2^{-193}$, the key-guessing
term satisfies $D \Nprim / 2^{256} \leq 2^{126}/2^{256} = 2^{-130}$, the
remaining linear terms satisfy
$(\Qv + \Nleaf)/2^{256} \leq 2^{64}/2^{256} = 2^{-192}$, and
$D(D-1)/2^{257} < 2^{-131}$; the total is below
$5 \cdot 2^{-131} + 2^{-192} + 2^{-193} < 2^{-128}$. These estimates
instantiate every closed form in \Cref{sec:concrete};
\Cref{sec:bound-discussion} examines the margin they leave.
\end{proof}

\begin{table}[t]
\centering
\caption{Concrete \TW{} security at the work cap $\Ntot \leq 2^{63}$
(with $D \leq 2^{63}$ for mu-ind), under
$k = c = \nu = \tauroot = \tauleaf = 256$ and the resource consequences
$\Nprim, \Qv, \Nleaf \leq \Ntot$ and $\Nprim + \Nleaf \leq \Ntot$ of
\Cref{sec:resources}. The Coll and
CMTDs-4 rows hold for every $s \geq 2$; larger $s$ only shrinks the
root-collision term.}
\label{tab:concrete-security}
\begin{tabular}{@{}lcc@{}}
\toprule
Notion & Dominant scale & At the cap \\
\midrule
All-in-one AE & $\Lift_c$ & $\leq 2^{-128}$ \\
Multi-user ind. & $D\,\KeyGuess_k(\Nprim)$ & $\leq 2^{-128}$ \\
\textsf{INT-RUP} & $\Lift_c$ & $\leq 2^{-128}$ \\
$\HTW$ collision (Coll) & $\Lift_c$, $\LeafColl$, $\RootColl$ & $\leq 2^{-128}$ \\
$\HTW$ second-preimage (Sec/eSec) & inherited from Coll & $\leq 2^{-128}$ \\
$\HTW$ preimage (Pre) & $\Lift_c$, $\LeafColl$ & $\leq 2^{-128}$ \\
$\HTW$ everywhere preimage (ePre) & $\Lift_c$ & $\leq 2^{-128}$ \\
CMTDs-4 \cite{BH22} & $\Lift_c$ and $\RootColl$ & $\leq 2^{-128}$ \\
\bottomrule
\end{tabular}
\end{table}

\subsection{Discussion of the Bounds}
\label{sec:bound-discussion}

The bounds of \Cref{sec:concrete} aggregate the loss terms of
\Cref{sec:loss-terms}, whose sizes at the work cap
\Cref{tab:concrete-security} records and whose dominance
\Cref{sec:concrete-claim} proves.

\paragraph{At the concrete cap.}
At $\Ntot \leq 2^{63}$, the hash bounds other than ePre and the
CMTDs-4 bound have flat-sponge
and root or leaf birthday terms on the same $2^{-131}$ scale; the
ePre and \textsf{INT-RUP} bounds retain the flat sponge loss
as their only birthday-scale term, and the single-user AE bound carries
its two flat-sponge terms, at $\Ntot$ and $\Nprim$, as its only
birthday-scale terms. Thus the flat sponge loss is not always the largest
summand: under the loose per-term cap, the $s=2$ root-collision term
$\binom{\Nprim+2}{2}/2^{256}$ is on the same scale as $\Lift_c(\Ntot)$.
The multi-user bound is the exception: its
hybrid-induced
term $D\,\KeyGuess_k(\Nprim)$ can reach $2^{-130}$ at the extreme cap
$D=\Nprim=2^{63}$. That corner is a property of the bound, not of any
admissible attack: the matching attack of the Tightness paragraph below
needs one accepted nonempty encryption per targeted user, and even the
empty-AD one-block case costs two duplex calls (root initialization and
closing root message). With $t$ targeted users and $q=\Nprim$ direct
guesses, the joint cap imposes $2t+q \leq 2^{63}$, so the feasible
key-guessing attack scale is about $2^{-133}$, reached near
$t=2^{61}$ and $q=2^{62}$. The remaining factor is hybrid slack from
charging created-but-idle users. Expressed in bytes, $\Ntot \leq 2^{63}$ corresponds to
at most $\rho \cdot 2^{63} = 167 \cdot 2^{63} < 2^{71}$ bytes of
plaintext, associated data, ciphertext, and aggregation through the
duplex (\Cref{sec:params}).

\paragraph{Margin at the cap.}
\Cref{cor:concrete-security} clears its target with a stated remainder:
at the joint corner $D = \Nprim = 2^{63}$ the multi-user estimate has
exact value $5 \cdot 2^{-131} + 2^{-194}$, since the corner forces
$N=0$ and hence $\Qv=\Nleaf=0$; this is just over five eighths of the
$2^{-128}$ budget. The largest of the remaining bounds, the $s = 2$
collision bound, stays below $4 \cdot 2^{-131}$
(\Cref{sec:concrete-claim}). The closeness at the corner is a proof
artifact rather than a consequence of the design. The corner itself is
unoccupiable: it forces $N = 0$, as the previous paragraph records, and
an adversary without construction computations receives identical
answers in both worlds---the primitive interface is the same real
permutation (\Cref{sec:games-pp})---so its advantage there is exactly
$0$. At resource splits an adversary can occupy, the multi-user
key-guessing witness reaches about the $2^{-133}$ scale, while the
root-collision and flat-sponge tightness witnesses can still reach the
$2^{-131}$ scale. The intervening distance to the proof bound is
consumed by the two flat-sponge composition terms and the hybrid slack
recorded above, the overhead of the indifferentiability route. The margin also widens away
from the corner: at $D \leq 2^{32}$, the $D$-indexed terms drop below
$2^{-160}$ and the multi-user bound falls to its two flat-sponge terms,
below $2^{-130} + 2^{-159}$, two bits under the target. A direct
keyed-duplex analysis would be expected to replace the quadratic
$\Lift_c$ terms by finer construction--primitive cross terms, much as
the keyed-duplex bounds synthesized by \cite{Men23} improve on the
sponge-indifferentiability bound underlying the original SpongeWrap
analysis of \cite{BDPV11}; \Cref{sec:limitations} discusses that route
and the gaps it would need to close.

\paragraph{Reading the work cap.}
\Cref{cor:concrete-security} is a success-probability statement at a fixed
cap, not a statement that security is exhausted there. The closed forms of
\Cref{sec:concrete} remain explicit functions of the resources beyond
$2^{63}$, and their dominant terms are birthday-type in $\Ntot$ on the
$256$-bit capacity, so the guarantees degrade quadratically rather than
abruptly: at $\Ntot = 2^{100}$, for example, the $s = 2$ collision bound is
still below $2^{-55}$. The bounds become vacuous only as $\Ntot$ approaches
$2^{c/2} = 2^{128}$, where $\Lift_c(\Ntot)$ exceeds $1/2$
(\Cref{sec:loss-terms}). The cap $\Ntot \leq 2^{63}$ is therefore a
data-scale choice: a single work bound at which every notion's advantage
lies below the $2^{-128}$ target simultaneously, not the largest workload
the mode tolerates.

\paragraph{Tightness.}
\begin{itemize}
\item \emph{$\KeyGuess_k(Q) = Q / 2^k$.} Tight at constants: a direct
  random-guess attack on a $k$-bit key matches the per-query success
  probability $1 / 2^k$. The multi-user bound's dominant term
  $D\,\KeyGuess_k(\Nprim) = D\Nprim / 2^k$ is likewise tight at the
  leading constant in the theorem's free-parameter setting, using the
  per-user nonce discipline of \Cref{sec:mu-ind}, which permits one
  nonce value across users. The adversary queries every targeted user
  on the same nonce $U$, empty associated data, and a known nonempty
  message---construction work charged to $N$, which does not enter the
  term---so each user's first ciphertext block reveals the keystream
  emitted by that user's root initialization call
  (\Cref{tab:transcript-schedule}). Each direct guess is then a single
  forward primitive query at the root initialization transcript
  $\prfxroot \concat K' \concat U$ for a fresh candidate key $K'$,
  whose output is compared against all $D$ collected keystream blocks
  at once; it matches some targeted key with probability $D / 2^k$, so
  $\Nprim$ guesses recover a user key---and distinguish---with
  probability $\approx D\Nprim / 2^k$, matching the bound. Under the
  joint cap the construction work for those target encryptions lowers
  the feasible scale to about $2^{-133}$, as noted in the concrete-cap
  paragraph above.
\item \emph{$\RootColl_\tau(Q, s) = \binom{Q + s}{s} / 2^{\tau (s -
  1)}$.} Tight up to lower-order terms: it is the standard $s$-way
  multi-collision bound on a $\tau$-bit projection of $\ROflat$, and
  a generic birthday-style adversary against the random oracle
  matches the leading-order term. It is the root term of the collision
  bound ($\HTW$ Coll, \Cref{thm:coll}). Its $\Nprim$ counts every
  primitive query, though by Output-Point Separation
  (\Cref{cor:output-sep}) only those whose induced $\ROflat$ call lands at
  a root tag output point ($\Rtag$, or $\Rmsg^+$ for single-chunk
  ciphertexts) can enter the candidate sequence; a class-filtered count
  would be no larger, and potentially smaller, but the displayed bound keeps the
  $\Nprim$ form for uniformity with \Cref{cor:ae,thm:mu-ind}, the birthday
  threshold at $\tauroot = 256$ being unaffected.
\item \emph{$\RootPre_\tau(Q) = (Q + 1) / 2^\tau$.} Tight at constants: a
  generic preimage search against a fixed $\tau$-bit target matches the
  per-query $2^{-\tau}$ rate. It is the root term of the standard and
  everywhere preimage bounds, where a single target tag yields a linear term.
\item \emph{Leaf tag term $\Nleaf / 2^{\tauleaf}$.} Tight at constants
  by a nonce-respecting attack: encrypt a single two-chunk message
  under nonce $U$, obtaining ciphertext with one leaf chunk and root
  tag $T$, then submit non-replay verification queries that each
  replace the leaf chunk by fresh bytes of the same length, keeping
  $T$. Each query recomputes the leaf tag at the hidden slot
  $(K, U, 1)$ and accepts whenever that fresh tag equals the
  encryption's hidden leaf tag $T_1$, since the aggregation
  string---and hence the recomputed root tag---is then unchanged. Each
  verification-side final leaf transcript therefore forges with
  probability $2^{-\tauleaf}$, matching the bound's per-transcript
  rate.
\item \emph{Known-key leaf collision $\LeafColl_\tau(Q)$.} Genuinely
  birthday-tight: when keys are adversary-supplied, direct queries can
  target leaf tag points, and a birthday search over chosen leaf bodies
  yields two distinct $\HTW$ inputs sharing a leaf tag and hence
  colliding openings; the collision proof uses $Q=\Nprim+\Nleaf$.
\item \emph{Flat sponge lift.} The term
  $\Lift_c(\Ntot) = \Ntot(\Ntot + 1) / 2^{c + 1}$ dominates the
  random-permutation differentiating bound $f_P(\Ntot)$ of
  \cite[Lemma~5]{BDPV08} via \Cref{lem:lift-bound}. The
  differentiating bound $f_P(\Ntot)$ is itself matched at leading order
  by the standard inner-collision differentiator, which succeeds
  against every simulator because the ideal-world construction
  interface is the random oracle regardless of simulator strategy; any
  residual conservatism in $\Lift_c(\Ntot)$ lies in composing
  indifferentiability into the specific keyed games, not in an
  improvable simulator bound, and is not specific to \TW{}. The
  real-vs-ideal games pay the same differentiating bound once more, at
  $\Nprim$, to exchange the ideal-side simulator for the real
  permutation (\Cref{lem:sim-perm-proximity}); since
  $\Nprim \leq \Ntot$, that term never dominates the $\Ntot$ term.
\end{itemize}

Cross-scheme bound-level comparisons against prior sponge AEADs and
the committing AEAD line of \cite{BH22} are deferred to
\Cref{sec:compare-table,sec:compare-sponge,sec:compare-commit}.
